%!TEX root = ../../thesis.tex

\section{Reinforcement Learning}\label{sec:reinforcement-learning}

% Drafting outline: bullet points to be fleshed out into prose.

\subsection{Agent, Environment, and Markov Decision Processes}\label{sec:rl-mdp}

% --- Verbatim excerpt from the proposal (own text; citation keys remapped) ---
Reinforcement learning (RL) is an extremely general framework for approaching problems
in which an agent strives to maximize a reward signal through a selection of actions
in certain states~\cite{sutton_reinforcement_2020}.
These kinds of problems can be modeled as finite Markov decision processes (MDPs).
Formally, a Markov decision process is a 3-tuple $\mathit{MDP}=(\mathcal{S},\mathcal{A},p)$ where

\begin{itemize}
  \item $\mathcal{S}$ is the set of states,
  \item $\mathcal{A}$ is the set of actions where $\mathcal{A}_s\subseteq\mathcal{A}$ are actions available in state $s\in\mathcal{S}$,
  \item $p(s^\prime,r|s,a)=\Pr\{S_t=s^\prime, R_t=r|S_{t-1}=s, A_{t-1}=a\}$
    describes the system dynamics, i.e., the probability of encountering a next
    state $s^\prime$ and reward $r$ at time step $t$ when taking action $a$ in
    state $s$ at time step $t-1$. This satisfies the \emph{Markov property}:
    $p$ depends only on the current state and not the full history.
\end{itemize}

An RL agent's goal is to learn a policy $\pi:\mathcal{S}\times\mathcal{A}\rightarrow [0,1]$
mapping each state $s$ to probabilities of taking each possible action $a$ to maximize
the cumulative reward.

A common way to visualize MDPs is via a directed graph: nodes model the state
space $\mathcal{S}$ while edges $(s_i,s_j)$ labeled with an action-reward-probability
tuple $(a,r,p)\in\mathcal{A}\times\mathcal{R}\times[0,1]$ model potential
transitions.

A standard example introduced in~\textcite{sutton_reinforcement_2020} is that of a \emph{recycling robot} (\cref{fig:mdp-recycling-robot}).
The robot collects empty cans and, at each decision point, either actively searches
for one, waits for someone to hand it a can, or returns to its charging station.
Its state is the charge level of its battery, $\mathcal{S}=\{\mathtt{high},\mathtt{low}\}$.

%!TEX root = ../thesis.tex
% Second example MDP for \cref{sec:rl-mdp}: the recycling robot of Sutton and Barto
% (Example 3.3), redrawn in the (a,r,p) edge convention of \cref{fig:mdp-example}.
\begin{figure}[htbp]
  \centering
  \begin{tikzpicture}[
      >={Stealth[round]},
      every state/.style={draw,minimum size=1.1cm,inner sep=1pt},
      every edge/.style={draw,->,font=\small},
      every loop/.style={->,looseness=8},
    ]
    \node[state] (high) at (0,0)   {$\mathtt{high}$};
    \node[state] (low)  at (7.6,0) {$\mathtt{low}$};

    \path
      (high) edge[loop above] node[above] {$(\mathtt{search},r_{\mathrm{s}},\alpha)$}      (high)
      (high) edge[loop below] node[below] {$(\mathtt{wait},r_{\mathrm{w}},1)$}             (high)
      (low)  edge[loop above] node[above] {$(\mathtt{search},r_{\mathrm{s}},\beta)$}       (low)
      (low)  edge[loop below] node[below] {$(\mathtt{wait},r_{\mathrm{w}},1)$}             (low)
      (high) edge[bend left=32] node[above] {$(\mathtt{search},r_{\mathrm{s}},1-\alpha)$}  (low)
      (low)  edge node[above] {$(\mathtt{recharge},0,1)$}                                  (high)
      (low)  edge[bend left=32] node[below] {$(\mathtt{search},-3,1-\beta)$}               (high);
  \end{tikzpicture}
  \caption[Example MDP: the recycling robot]{%
    The recycling robot of Sutton and Barto~\cite{sutton_reinforcement_2020}.
    Its state space consists of the charge level of its battery, $\mathcal{S}=\{\mathtt{high},\mathtt{low}\}$,
    and its action space consists of $\mathcal{A}=\{\mathtt{search},\mathtt{wait},\mathtt{recharge}\}$.
    Its battery state transition is probabilistic and depends on $\alpha$ and $\beta$.
    $r_s$ and $r_w$ are the rewards for search and waiting.
  }
  \label{fig:mdp-recycling-robot}
\end{figure}


A simple greedy approach is to always select the $\mathtt{search}$ action since it
pays the highest rewards in expectation.
This approach incurs the penalty of searching while the battery is in the $\mathtt{low}$ state.
An agent must therefore recognize the state it is in and choose its action accordingly.
While this example is easily solved, MDPs with larger state and action spaces present
a much harder problem.

\subsubsection{Partially Observable Markov Decision Processes}\label{sec:rl-pomdp}

In the previous formulation of MDPs, the true state of the environment is available
to the agent.
In reality, this is often not the case.
An agent perceives its environment as a lossy function of the underlying state.
The recycling robot of \cref{fig:mdp-recycling-robot}, for instance, might perceive
its charge level as a noisy voltage measurement.
Observing the environment may therefore introduce uncertainty.
The Markov property then still holds for the underlying state, but not for what the agent sees.

This situation is modeled through a \emph{partially observable Markov decision process}
(POMDP)~\cite{ASTROM1965174,kaelbling_planning_1998}, which
extends the original 3-tuple MDP $(\mathcal{S},\mathcal{A},p)$ to a 5-tuple $\mathit{POMDP}=(\mathcal{S},\mathcal{A},p,\mathcal{O},o)$, by adding

\begin{itemize}
  \item $\mathcal{O}$, the set of observations the agent can receive, and
  \item $o(\omega|s^\prime,a)=\Pr\{O_t=\omega|S_t=s^\prime,A_{t-1}=a\}$, the observation
    function, giving the probability of perceiving $\omega$ after a transition into
    $s^\prime$ under action $a$.
\end{itemize}

The environment still evolves as an MDP over $\mathcal{S}$ according to $p$.
However, the agent must act on observation $O_t$ instead of on the true state $S_t$.
An ordinary MDP is the special case $\mathcal{O}=\mathcal{S}$ with $o=\operatorname{id}(\cdot)$.

POMDPs introduce the problem of \emph{state aliasing}.
If two states that require different actions map to the same observation, then no deterministic
memoryless policy $\pi:\mathcal{O}\times\mathcal{A}\rightarrow[0,1]$ can act correctly in both,
and the best achievable memoryless policy can be arbitrarily worse than a policy that is
allowed to consult the history~\cite{SINGH1994284}.

A POMDP can be made Markovian by treating not the actual observations, but the
entire observation-action history $H_t=O_0,A_0,O_1,A_1,\dots,O_t$ as the state space.
Another classical solution is to summarize the observation-action history in a probability distribution
and treat this distribution as the actual state.
This state is called a \emph{belief state} $b_t(s)=\Pr\{S_t=s|H_t\}$ and suffices
to allow the agent to act optimally~\cite{puterman_markov_2009}.
Practically speaking, the belief-state approach is restricted in that it presupposes
the existence of a known model $(p,o)$. Solving the resulting belief MDP is also intractable
for all but the simplest problems~\cite{papadimitriou_complexity_1987}.

\subsection{The Exploration-Exploitation Tradeoff}\label{sec:rl-exploration}

The most important problem any RL agent encounters is that of the \emph{exploration-exploitation tradeoff}.
At any point in time, the agent can either \emph{exploit} previously gained information
to achieve a higher reward, or \emph{explore} yet unknown actions to gain more information
that might come in useful in the future.
Balancing the selection of actions that explore vs. actions that exploit is a key tradeoff in any RL algorithm.

Any MDP where $\lvert\mathcal{S}\rvert=1$ and $\lvert\mathcal{A}\rvert=k$ is also known as
a \emph{$k$-armed bandit}, where the $k$-th arm is equal to the $k$-th action $a_k\in\mathcal{A}$.
This view is inspired by a slot machine with $k$ arms, where pulling a lever produces a reward $r_k\sim R_k$
drawn from a stationary distribution $R_k$.
In this setting, the difficulty lies in judging which arm is best.
Each arm has an expected value $q^*(a_k)=\mathbb{E}[R_t|A_t=a_k]$ that is unknown
and must be estimated, for example by the sample average

\begin{equation}\label{eq:bandit-sample-average}
  Q_t(a_k)=\frac{1}{N_t(a_k)}\sum_{\substack{i<t\\A_i=a_k}}R_i,
\end{equation}

where $N_t(a)$ denotes the number of times arm $a$ has been pulled before time step $t$.
The exploration-exploitation tradeoff in this case manifests itself through the choice
of whether more samples should be collected to improve the estimate or whether
the current estimate should be used to select actions to gain more reward.
We will outline three possible strategies.

\paragraph{Greedy}

A greedy agent always pulls the arm with the highest current estimate

\[
  a_t^*=\argmax_{a} Q_t(a)
\]

This strategy can easily fail.
If the agent is unlucky and sees a few bad initial rewards, the sample average $Q_t$ will be off and $a_t^*$ will be suboptimal.

\paragraph{$\varepsilon$-greedy}
The simplest remedy is to explore at random.
With probability $\varepsilon$, the agent pulls an arm chosen uniformly at random, and exploits greedily otherwise.
This guarantees that every arm keeps being sampled, but the exploration is undirected.
Moreover, even after the estimate has converged, the agent keeps exploring.
This is typically countered by decaying $\varepsilon$ over time~\cite{sutton_reinforcement_2020}.

\paragraph{Upper confidence bounds}
Directed exploration makes use of the \emph{upper confidence bounds} rule to estimate which arms
could still yield a high reward given how little it has been tried.
The rule is formulated in terms of the current estimate plus the visit count of each action
and prioritizes high-value actions with low visit count.

\begin{equation}\label{eq:ucb}
  A_t=\argmax_a\left[Q_t(a)+c\sqrt{\frac{\ln t}{N_t(a)}}\,\right],
\end{equation}

This rule has been proven to be optimal in terms of regret minimization up to a constant factor~\cite{auer_finite-time_nodate}.
UCB is generalizable from bandits to full MDPs by treating every decision point
of a search tree over the action space as its own bandit.
This gives rise to the UCT rule~\cite{hutchison_bandit_2006} that is used in an
algorithm called Monte-Carlo Tree Search (MCTS)~\cite{hutchison_efficient_2007}.

\subsection{Solving MDPs}\label{sec:rl-solving}

Methods for solving MDPs can be classified along two axes.

\paragraph{Representation of estimates}
Classical algorithms represent estimates via tables.
The size of the tables usually scales with action and state space size.
This makes \emph{tabular methods} applicable only to small problems with tractable size.
\emph{Function approximation methods} exploit regularities in the table data by fitting a parameterized model.
Earlier work that made use of manually engineered features has been succeeded by modern, deep-learning based approaches.

\paragraph{Model availability}
Some algorithms require that a model of the transition function $p$ is available.
Classically, the model is assumed to be specified. In more modern algorithms,
the model is learned by the agent itself through interactions with the underlying MDP.
\emph{Model-free methods} don't base their estimates on an underlying model, but
derive them directly from the interactions.

%!TEX root = ../thesis.tex

\begin{table}[htb]
  \centering
  \caption[Classification of \gls{rl} algorithms]{Classification of \gls{rl} algorithms
    along the model axis and the representation axis, with representative
    algorithms.}
  \label{tab:rl-taxonomy}
  \begin{tabularx}{\textwidth}{@{}lXX@{}}
    \toprule
    & Model-based & Model-free \\
    \midrule
    Tabular
      & Policy/value iteration\newline(\cref{sec:rl-dp})
      & \gls{sarsa}, Q-learning\newline(\cref{sec:rl-td}) \\
    \addlinespace
    Function approximation
      & MuZero, EfficientZero\newline(\cref{sec:rl-sample-efficiency})
      & \gls{dqn}, \gls{ppo}\newline(\cref{sec:rl-fa}) \\
    \bottomrule
  \end{tabularx}
\end{table}


\subsubsection{Value functions and Bellman equations}\label{sec:rl-value-functions}

We will introduce some definitions and formalize what it means for a policy to be optimal.

The \emph{return} at time step $t$ is the discounted sum of all future rewards,

\begin{equation}\label{eq:return}
  G_t=R_{t+1}+\gamma R_{t+2}+\gamma^2 R_{t+3}+\dots=\sum_{k=0}^{\infty}\gamma^k R_{t+k+1},
\end{equation}

where the \emph{discount factor} $\gamma\in[0,1)$ controls how far-sighted the agent is
by exponentially discounting future rewards.

A \emph{value function} measures how good a state is under a given policy.
It is defined as the expected return when starting in $s$ and following $\pi$:

\begin{equation}\label{eq:value-function}
  v_\pi(s)=\mathbb{E}_\pi[G_t|S_t=s]
\end{equation}

The \emph{Bellman equation} is a recursive consistency equation that formalizes the
intuition of the value of a state being equal to its expected immediate reward plus
the discounted value of its successor states:

\begin{equation}\label{eq:bellman-v}
  v_\pi(s)=\sum_a\pi(a|s)\sum_{s^\prime,r}p(s^\prime,r|s,a)\left[r+\gamma v_\pi(s^\prime)\right].
\end{equation}

An \emph{optimal policy} $\pi^*$ has the property that no other policy assigns a
higher value to any state: $\forall\pi\;\forall s\in\mathcal{S}:v_{\pi^*}(s)\geq v_\pi(s)$.
All optimal policies have the same value function, denoted as $v^*=v_{\pi^*}$.

The recursive \emph{Bellman optimality equation} follows:

\begin{equation}\label{eq:bellman-opt}
  v^*(s)=\max_a\sum_{s^\prime,r}p(s^\prime,r|s,a)\left[r+\gamma v^*(s^\prime)\right],
\end{equation}

We can define the same equations over state-action pairs $(s,a)\in\mathcal{S}\times\mathcal{A}$ instead:

\begin{align}
  q_\pi(s,a)&=\mathbb{E}_\pi[G_t|S_t=s,A_t=a]\label{eq:q-function}\\
  q_\pi(s,a)&=\sum_{s^\prime,r}p(s^\prime,r|s,a)\left[r+\gamma\sum_{a^\prime}\pi(a^\prime|s^\prime)\,q_\pi(s^\prime,a^\prime)\right]\label{eq:bellman-q}\\
  q^*(s,a)&=\sum_{s^\prime,r}p(s^\prime,r|s,a)\left[r+\gamma\max_{a^\prime}q^*(s^\prime,a^\prime)\right]\label{eq:bellman-opt-q}
\end{align}

Given $v^*$ or $q^*$, we can solve for $\pi^*$ by either choosing the action that leads
to the highest-valued state or directly choosing the highest-valued action.

\subsubsection{Dynamic programming}\label{sec:rl-dp}

When the model $p$ is known, the Bellman optimality equation can be used to derive an algorithm
that provably converges to $\pi^*$.

\paragraph{Policy iteration}
Let $\pi_i$ be given. By repeatedly applying the Bellman equation,
we can compute $v_{\pi_i}$. Once $v_{\pi_i}$ is known, we improve $\pi_i$ by acting greedy
with respect to $v_{\pi_i}$. This then yields a strictly better policy $\pi_{i+1}$.
A fixpoint after $n$ steps is then guaranteed to be optimal: $\pi_n=\pi^*$.

\paragraph{Value iteration}
Instead of alternating between deriving $v_{\pi_i}$ and improving $\pi_i$, we can directly
apply the Bellman optimality equation to get a one-step update rule~\cite{puterman_markov_2009}:

\begin{equation}\label{eq:value-iteration}
  v_{k+1}(s)=\max_a\sum_{s^\prime,r}p(s^\prime,r|s,a)\left[r+\gamma v_k(s^\prime)\right].
\end{equation}

Most RL algorithms make use of this general framework of estimating values and then
using these estimates to improve the current policy.
This framework is called \emph{generalized policy iteration} (GPI)~\cite{sutton_reinforcement_2020}.

While Dynamic Programming guarantees convergence to an optimal policy, it requires a complete model $p$ and
must sweep the entire state space each iteration, making it impractical for larger real-world problems.

\subsubsection{Temporal-difference learning}\label{sec:rl-td}

Temporal-difference (TD) learning removes the need for a model.
Values are estimated from sampled transitions $(S_t,A_t,R_{t+1},S_{t+1})$ obtained by interacting with the environment~\cite{sutton_learning_1988,sutton_reinforcement_2020}.
TD methods, like DP, make use of \emph{bootstrapping} --- updating estimates based on other estimates.
They don't require full state space sweeps, but update the value function using the sampled transition data.
A simple TD algorithm, TD(0), updates the estimate $V(S_t)$ as follows:

\begin{equation}\label{eq:td0}
  V(S_t)\leftarrow V(S_t)+\alpha\,\delta_t,
  \qquad
  \delta_t=R_{t+1}+\gamma V(S_{t+1})-V(S_t),
\end{equation}

where $\alpha\in(0,1]$ is a step size and the \emph{TD error} $\delta_t$ measures how much the
sampled outcome deviates from the current estimate.

Monte Carlo methods form the bootstrap-free alternative by simulating a whole
episode, updating toward the full observed return $G_t$ at the end.
The TD($\lambda$) family interpolates between the two extremes via \emph{eligibility traces},
with TD(0) and Monte Carlo as the limit cases $\lambda=0$ and
$\lambda=1$~\cite{sutton_learning_1988}.

\paragraph{SARSA}
SARSA~\cite{sutton_emphatic_2015} is similar to TD learning, but learns an action-value function $Q(s,a)$ instead of
a value function $V(s)$. The update rule makes use of the action that the agent
actually performed, making SARSA an \emph{on-policy} algorithm since it follows
the currently active policy $\pi$:

\begin{equation}\label{eq:sarsa}
  Q(S_t,A_t)\leftarrow Q(S_t,A_t)+\alpha\left[R_{t+1}+\gamma Q(S_{t+1},A_{t+1})-Q(S_t,A_t)\right].
\end{equation}

\paragraph{Q-learning}
Q-learning\cite{watkins_q-learning_1992} can be interpreted as an \emph{off-policy} SARSA:
its target maximizes over the successor actions, making the update rule
independent from the actually taken action:

\begin{equation}\label{eq:q-learning}
  Q(S_t,A_t)\leftarrow Q(S_t,A_t)+\alpha\left[R_{t+1}+\gamma\max_{a}Q(S_{t+1},a)-Q(S_t,A_t)\right],
\end{equation}

\subsubsection{Function approximation}\label{sec:rl-fa}

The three algorithms presented above are tabular: they maintain one entry per state or state-action pair,
requiring a state space small enough to enumerate.
Realistic state spaces are far too large for tables.
Additionally, tabular methods cannot generalize to new, yet unseen states.
Function approximation addresses both problems by replacing the table with a parameterized
function $\hat{v}(s;\theta)$ or $\hat{q}(s,a;\theta)$ whose parameter vector $\theta$
can be learned, for example by (semi-)gradient descent on the TD error.

Historically, the approximators were linear in hand-crafted features.
Today, they are deep neural networks, giving rise to the field of \emph{deep reinforcement learning}.
The combination is powerful, but introduces new issues concerning stability:
bootstrapping, off-policy learning, and function approximation together can
make training diverge~\cite{sutton_reinforcement_2020}.

The two large families of model-free deep reinforcement learning algorithms are
the deep Q-network (DQN)~\cite{mnih_human-level_2015} and policy gradient
methods such as proximal policy optimization
(PPO)~\cite{schulman_proximal_2017}. DQN uses the standard Q-learning update
rule, but represents $Q$ as a function approximator instead, allowing it to
deal with large state and action spaces. PPO directly learns a parameterized
policy $\pi_\theta$ by performing gradient updates on a clipped objective
function that increases the probability of advantageous actions while limiting
the magnitude of policy changes.

\subsubsection{Model-based deep reinforcement learning and sample efficiency}\label{sec:rl-sample-efficiency}

Model-free deep RL algorithms require a large amount of data.
For instance, DQN and its successors require on the order of tens to hundreds
of millions of environment interactions to learn a useful $Q$ function.
This is a major drawback if environment interactions are expensive.
Many real-world applications therefore require high \emph{sample efficiency}.

Model-based deep RL algorithms allow higher sample efficiency by learning a model of the underlying dynamics
while they interact with the environment.
This allows them to generate additional, cheap experience~\cite{moerland_model-based_2022},
and to train the policy inside the learned environment.
The Dreamer line refines this recipe of learning latent
dynamics and imagining rollouts~\cite{hafner_learning_2019,hafner_mastering_2024}.
The MuZero line instead learns an implicit latent model end-to-end and uses it for planning~\cite{schrittwieser_mastering_2020}.
Both approaches achieve remarkable results with very few samples, making them
great candidates for when sample efficiency is a high priority.
